\worldchapter{\trans{recipes_title}}{brewing}
\label{ch:appendix-recipes}

\begin{multicols}{2}


    \section{Potions and Tinctures}



    \subsection*{Laughing Potion}

    \begin{tabular}{@{}ll@{}}
        \textbf{\trans{difficulty_label}:} & simple \\
        \textbf{\trans{yield_label}:} & 1 \\
    \end{tabular}

    Preparation:
First, bring the spring water to a boil in a copper kettle until it bubbles gently. Add the crushed comfrey, which forms the physical basis and ensures that the laugh muscles remain supple.

Now comes the crucial moment: the “giggle pea” must be crushed in a mortar while telling a (really!) bad joke out loud. The flatter the joke, the more sparkling the potion! Stir the mixture three times counterclockwise while giggling quietly. Finally, stir in the honey to dispel the bitterness of the world.

Preparation time: 30 minutes of active stirring and joke telling, followed by 1 hour of resting time to “settle.”

Quantity: 1 application, approx. 50 Oth liquid, weight: 50 Oth.

Effect:
The consumer must laugh loudly at every little incident at the next opportunity (duration: 2D6 minutes). In terms of game mechanics, this grants a bonus die on all “Performance” or ‘Communication’ checks, but makes all “Stealth” checks +3 more difficult (SG Medium).

Exhaustion: After the effect expires, -1 point on all dice rolls for 30 minutes (laughing is exhausting!).

    \paragraph{\trans{ingredients_label}}
    \begin{denseitemize}
        \item \textbf{Water} (60 oth), \textit{Components}.
        Cold, clear water.
        \item \textbf{Comfrey (Symphytum officinale)} (10 oth), \textit{Components}.
        Comfrey stimulates blood circulation, bruises, hematomas and sprains disappear faster. Comfrey accelerates the regeneration of cells.
        \item \textbf{Chickpeas (Gigglepea)} (1 Piece), \textit{Food / Provisions}.
        This small, yellow-green bean is a truly wonderful food.

Apart from its alchemical use as a base for laughing potions, it is an excellent food for weary travelers.

It is often found in the rich markets in and around Al Bah Ji Ra. However, its popularity is slowly spreading throughout Tirakan.

Eating a handful of roasted chickpeas dispels gloomy thoughts and briefly lifts the “shadow” from the mind. This is especially true when seasoned with spices such as garlic, cumin, or pepper.

Even on their own, they are rich in protein and complex fiber, making them a nutritious addition to any diet.

“If the stew tastes too bitter and life weighs too heavily, throw in a handful of chickpeas. Your stomach will thank you with a gurgle.”
— Master chef Alar al-Din of the “Tent of Seven Veils” in El Kurru
        \item \textbf{Honey} (5 oth), \textit{Food / Provisions}.
        In cultivated regions, honey is produced by hard-working beekeepers, who often set up their hives near herb gardens.

This honey is clear, stable, and has a delicate taste of thyme or lavender. It usually costs about 2 guilders per 5 oth and is a reliable staple for any alchemist who wants to soften the bitterness of herbal extracts.

Those who are more adventurous search the forests for wild bee nests. This honey is darker, thicker, and often mixed with pollen or small pieces of wax. It has a strong, almost earthy flavor. It is said that wild honey from the primeval forests has a stronger regenerative power for the voice. Ideal for bards who need to lubricate their throats after a long night in the tavern.
    \end{denseitemize}


    \subsection*{Kinstarchel Secretion}

    \begin{tabular}{@{}ll@{}}
        \textbf{\trans{difficulty_label}:} & medium \\
        \textbf{\trans{yield_label}:} & 1 \\
    \end{tabular}

    Preparation:
First, the bones of a freshly deceased (or at least not too weathered) Kinstarchel must be carefully cleaned and crushed to extract the deep black marrow. This is heated over a low flame in pure alcohol. It must not be allowed to boil! As soon as an oily, slightly shimmering layer forms on the surface, it must be skimmed off with a silver pipette. The secretion is then mixed directly with a ready-made potion or poison. The mixture is highly unstable; a violent shock (such as the impact of a vial) is enough to suddenly release the energy it contains.

Preparation time: 1,5 hours of active extraction, after which the mixture must rest for 1 hour in a cool, dark container before it is stable enough for transport.

Quantity: 1 application (approx. 10 Oth secretion), sufficient to prepare one vial (50-100 Oth) of another potion.

Effect:
The secretion gives a potion the property “Explosive.” When the container is broken (Dexterity check required), the mixture explodes within a radius of 1D3 meters. Every creature in this area is treated as if it had directly ingested the potion or poison.

Note:
The secretion itself does not cause any damage; it merely serves as a rapid carrier for the effect of the mixed potion.

    \paragraph{\trans{ingredients_label}}
    \begin{denseitemize}
        \item \textbf{Alcohol} (300 oth), \textit{Components}.
        In the laboratories of alchemists and the huts of herbalists, alcohol is rarely stored for pleasure. It is considered a solvent capable of unleashing the essence of plants and minerals.

It is obtained from fermented grain or fruit through multiple distillations. It is a distillate with a pungent, sharp smell that is so strong that it tickles the nose when inhaled and evaporates immediately on the skin, leaving a cooling sensation.
        \item \textbf{Kinstarchel Bone} (2 stone), \textit{Components}.
        Bones of a deceased Kinstarchel. An explosive secretion can be extracted from the marrow of the bones.
    \end{denseitemize}


    \subsection*{Potion of Deep Calm}

    \begin{tabular}{@{}ll@{}}
        \textbf{\trans{difficulty_label}:} & medium \\
        \textbf{\trans{yield_label}:} & 2 \\
    \end{tabular}

    Preparation:
To brew this potion, first bring the water to a boil, then add the valerian and crushed moonthorn berries and boil for five minutes. Remove the mixture from the heat and let it steep for thirty minutes before bottling.

Perparation Time: 50 minutes, 6 hours resting time

Quantity: 2 Applications, 2x250 oth liquid, 2x250 oth weight

Effect:
The essence causes the character to fall into a deep, ten-hour sleep within five minutes, during which time D6 wounds are regenerated. To awaken from this deep slumber, the character must make two successful rolls of Resistance in a row when disturbed or attempts are made to wake them. Only direct physical force (shaking/attack) or extremely loud noises will cause immediate awakening.

    \paragraph{\trans{ingredients_label}}
    \begin{denseitemize}
        \item \textbf{Moonthorn Berry} (5 Berries), \textit{Components}.
        The moonthorn berry is a gift from the deepest night. It grows as a ground-covering shrub whose delicate tendrils and deep green leaves are protected by striking, short thorns.

It is found exclusively in places that rarely see the light of the sun, usually deep in ancient forests or near damp grotto and cave entrances. It only reveals its true splendor under the light of the full moon, when its small, berry-like fruits glow in a mysterious, dull blue, almost as if they had swallowed the light of the celestial sphere itself.

The berry is notorious for its strong sedative effect. In small doses, it has a calming effect, but when concentrated in a potion, its essence can numb the mind and put the body into a state of deep, dreamless stillness. Gathering them is risky, as their thorns can cause temporary itching when touched.
        \item \textbf{Valeriana (Valeriana officinalis)} (1 Bunch), \textit{Components}.
        Valerian helps with insomnia and restlessness. Hops and lemon balm increase the effect of valerian and improve the taste.
        \item \textbf{Water} (500 oth), \textit{Components}.
        Cold, clear water.
    \end{denseitemize}


    \subsection*{Simple wound tincture}

    \begin{tabular}{@{}ll@{}}
        \textbf{\trans{difficulty_label}:} & simple \\
        \textbf{\trans{yield_label}:} & 1 \\
    \end{tabular}

    Preparation: First, carefully clean the elecampane and comfrey. Finely slice the elecampane roots and crush them into a paste in a mortar, while the comfrey leaves should only be lightly rubbed between your fingers to release their juices. Place both in a ceramic bottle and fill it with clear spring water (or, if you want it to be particularly pure, with distilled spirit). Simmer the mixture over a gentle heat until the liquid takes on a deep green, almost oily consistency. Be careful! If you stoke the fire too hot, you will burn the healing spirits of the plants!

Preparation time:
Active time: 30 minutes (chopping and simmering). Resting time: 2 hours in a cool, dark place so that the active ingredients can be transferred to the tincture.

Quantity:
1 application, approx. 10 Oth liquid in a bottle.
Total weight: 10 Oth (including container).

Effect:
When used successfully with First Aid and a bandage, the bandage heals an additional 1W3 wounds.

    \paragraph{\trans{ingredients_label}}
    \begin{denseitemize}
        \item \textbf{Water} (100 oth), \textit{Components}.
        Cold, clear water.
        \item \textbf{Inula (Inula helenium)} (10 oth), \textit{Components}.
        This medicinal plant from the Middle Ages is no longer widely used in modern times. Its application improves digestion, and it is believed to have a preventive effect against colon cancer.
        \item \textbf{Comfrey (Symphytum officinale)} (10 oth), \textit{Components}.
        Comfrey stimulates blood circulation, bruises, hematomas and sprains disappear faster. Comfrey accelerates the regeneration of cells.
    \end{denseitemize}


    \subsection*{Magic potion (a full carafe)}

    \begin{tabular}{@{}ll@{}}
        \textbf{\trans{difficulty_label}:} & simple \\
        \textbf{\trans{yield_label}:} & 3 \\
    \end{tabular}

    Preparation:
First, prepare the brew from elecampane, comfrey, and yarrow. When the steam begins to burn heavily and herbaceously in your nose, add the scales of a river dragonfly. Stir seven times against the sun's course while thinking of cold, deep water. The potion will not glow like a cheap fairground trick, but will take on a deep, earthy blue that shimmers at the right angle.

Preparation time:
Active time: 30 minutes (scraping and cleaning the scales takes time)
Resting time: 2 hours 

Quantity:
5 applications, 500 Oth liquid, weight: 500 Oth

Effect: Restores 2 Arcana per application.

    \paragraph{\trans{ingredients_label}}
    \begin{denseitemize}
        \item \textbf{Inula (Inula helenium)} (10 oth), \textit{Components}.
        This medicinal plant from the Middle Ages is no longer widely used in modern times. Its application improves digestion, and it is believed to have a preventive effect against colon cancer.
        \item \textbf{Comfrey (Symphytum officinale)} (15 oth), \textit{Components}.
        Comfrey stimulates blood circulation, bruises, hematomas and sprains disappear faster. Comfrey accelerates the regeneration of cells.
        \item \textbf{Yarrows (Achillea millefolium)} (5 oth), \textit{Components}.
        Yarrow is used for its hemostatic effect. The flowers and the leaves contain tannins, bitter and mineral substances. The essential oil of the plant has anti-inflammatory and antispasmodic effect.
        \item \textbf{Scale of a river nymph} (1 Piece), \textit{Components}.
        The river nymph is an extremely shy creature, which, according to rumors in alchemical circles, no mortal being in Tirakan has ever truly seen. The only tangible proof of its existence are its scales.

They are occasionally found on muddy riverbanks, on rocks washed by spray, or in the darkness of underground lakes.

At dusk, these scales often appear simply greenish-brown, almost like ordinary horn or dried leaves. But as soon as the sun of Tirakan kisses the surface, they come to life and shimmer in all the colors of mother-of-pearl. They often reach the size of a proud palm and resemble the scales of a large fish in texture.

Fortunately, this magical ingredient is not a rare sight in the markets of the empire, especially on the trade routes between Toran and Yavon down to Meridian. Since they are found regularly, their value remains manageable, making them an honest ingredient. They are so common that it is hardly worthwhile for counterfeiters to produce inferior copies. But a word of warning from me: always look for the characteristic shimmer! 

According to an old legend, on lonely nights, river nymphs sit enthroned on mossy rocks and comb their endless, water-colored hair with combs made of bone and old driftwood. Their faces are said to be marked by a peaceful melancholy as they hum quiet songs whose deeper meaning is known only to the flowing water.
        \item \textbf{Water} (500 oth), \textit{Components}.
        Cold, clear water.
    \end{denseitemize}


    \subsection*{Guild Elixir}

    \begin{tabular}{@{}ll@{}}
        \textbf{\trans{difficulty_label}:} & simple \\
        \textbf{\trans{yield_label}:} & 3 \\
    \end{tabular}

    Preparation:
First, heat the pure spring water—do not bring it to a boil, just let it simmer! Finely chop the waxwort and add it to the hot water along with the honey.

Let the mixture steep until the honey has completely dissolved and the water takes on the golden color
of the herbs. Then strain the mixture through a fine linen cloth into the vial.

Preparation time:
1 hour of active brewing at the cauldron,
followed by half an hour of cooling time.

Quantity:
3 doses. 200 Oth of liquid, total weight 250 Oth (including glass vial).

Effect:
Taking one dose immediately alleviates existing exhaustion. Additionally, the user gains a +1 die bonus for the next 1W6 hours on all checks that require concentration or alertness
.

    \paragraph{\trans{ingredients_label}}
    \begin{denseitemize}
        \item \textbf{Honey} (20 oth), \textit{Food / Provisions}.
        In cultivated regions, honey is produced by hard-working beekeepers, who often set up their hives near herb gardens.

This honey is clear, stable, and has a delicate taste of thyme or lavender. It usually costs about 2 guilders per 5 oth and is a reliable staple for any alchemist who wants to soften the bitterness of herbal extracts.

Those who are more adventurous search the forests for wild bee nests. This honey is darker, thicker, and often mixed with pollen or small pieces of wax. It has a strong, almost earthy flavor. It is said that wild honey from the primeval forests has a stronger regenerative power for the voice. Ideal for bards who need to lubricate their throats after a long night in the tavern.
        \item \textbf{Water} (130 oth), \textit{Components}.
        Cold, clear water.
        \item \textbf{Gildenkraut (Vigilia mercatoris)} (1 Bunch), \textit{Components}.
        Gildherb is a tough, almost inconspicuous plant that only catches the trained eye of an expert.

It is characterized by jagged, leathery leaves of a deep green color, whose thick veins have a slight reddish sheen. Its flowers are tiny, pale yellow, and close immediately once the sun passes its zenith.

The herb is entirely mundane and possesses not the slightest spark of the arcane. Its value lies rather in its extremely high, natural concentration of alkaloids, particularly pure caffeine. When the tough leaves are crushed between the fingers, they give off a sharp, earthy scent reminiscent of heavily roasted nuts.

Location and Habitat:
This plant is relatively rare, as it requires very specific, adverse conditions to produce its stimulating sap. It thrives exclusively in barren, mineral-rich soils and avoids the damp moss of deep forests.

Gildherb is found primarily in the rocky crevices of ancient trade routes, in the crumbling stone mortar of abandoned guild outposts, or on the steep, sun-baked slopes of the Toran Highlands. It needs hard rock and plenty of wind to prevent it from withering away.
    \end{denseitemize}


    \subsection*{Major healing potion}

    \begin{tabular}{@{}ll@{}}
        \textbf{\trans{difficulty_label}:} & difficult \\
        \textbf{\trans{yield_label}:} & 3 \\
    \end{tabular}

    Preparation:
First, the pure spring water must be slowly brought to a boil over a fire made from dried oak wood. The dried sunblossom petals are finely ground in a mortar until they form a fine powder, which is then carefully stirred in. Next, the crushed bloodwort is added, which gives the potion its characteristic deep red color. The most critical moment is the addition of the tanium dust. The magical component must be added in a figure-eight motion while stirring constantly to bind the healing energies. Finally, the distillate is filtered through a fine silk cloth and poured into a vial. 

Preparation time: Active time: 3 hours (constant stirring and temperature control).

Resting time: 12 hours in complete darkness.

Quantity: 3 applications (enough for 3 separate healings). 

Liquid quantity: 15 oth per application (total 45 oth). 

Weight: approx. 50 oth. 

Effect: Healing: When applied, 1W3+3 wounds are healed.

    \paragraph{\trans{ingredients_label}}
    \begin{denseitemize}
        \item \textbf{Common Bloodweed} (10 oth), \textit{Components}.
        Bloodweed is a ground-level plant that is particularly striking due to its fleshy, deep red leaves.

The fine veins on the leaf surface glow in a rich scarlet red, almost as if real blood were pulsing through them. When a leaf is crushed, a sticky, sweet-smelling sap emerges that stains the fingers for days. It is not a magical plant in the classic sense. It draws its power from the iron-rich soil and the pale light of the dense forests.

You will usually find bloodweed in shady, damp places. It prefers the foot of old oak trees or the immediate vicinity of rotting undergrowth in deep forests. An inattentive traveler often mistakes it for common purple sorrel, but a trained alchemist will notice the small, pearl-like dewdrops that always collect at the edges of the leaves.

It is best harvested in the early morning hours, before the sun breaks through the canopy. Only the outer leaves are cut to preserve the root.
        \item \textbf{Taniumdust} (1 oth), \textit{Components}.
        Tanium is a dark, crystalline element of exceptional hardness. In its raw state, it is often found as deep black veins in ancient rock, mostly in areas with high concentrations of natural magic. 

It is so hard that conventional mortars break when used on it; only tools made of hardened diamond or magically reinforced grinding mechanisms can grind it into fine dust.

Tanium acts as a perfect storage medium for arcane energy. But beware: it has no saturation point! If it becomes saturated with too much magic or unstable due to impure alchemy, it will discharge in a magical explosion.
        \item \textbf{Sunblossom} (5 oth), \textit{Components}.
        The sunblossom is the epitome of constancy. With its strong, rough stem and proud, golden-yellow crown of petals, it tirelessly follows the course of the sun's chariot across the firmament of Tiraka.

Its core is filled with nutritious, oil-rich seeds, but for alchemists, it is the bright outer petals that are most valuable. They store the pure, gentle warmth of the day without carrying the dangerous heat of fire or the unpredictability of magic.

It can be found almost everywhere in the fertile plains of Tiraka, especially on the sun-drenched hills around Asgoran or in the gardens of farmers in the hinterland. It loves open spaces and deep, black soil. It is not a rare plant, no, but one that needs care. The wild varieties in the heaths are often smaller, but their essence is more concentrated than that of the cultivated specimens.

The leaves should be picked at midday, when the sun is at its highest and the flower is in full bloom. They should be dried flat on linen cloths in an airy place, never in direct oven heat!
        \item \textbf{Water} (45 oth), \textit{Components}.
        Cold, clear water.
    \end{denseitemize}


    \subsection*{Sud of shallow empowerment}

    \begin{tabular}{@{}ll@{}}
        \textbf{\trans{difficulty_label}:} & simple \\
        \textbf{\trans{yield_label}:} & 2 \\
    \end{tabular}

    Preparation:
The schnapps is mixed with the frost lichen until it takes on a bluish color. The blood of the flying lizard is only stirred in after the discoloration. The mixture must steep for at least 5 hours in a cool place (cellar, cave). Then pour it into the vial of flying lizard blood for the best effect!

Preparation time: 2 hours

Quantity: 2 applications, 2x200 Oth liquid, 2x200 Oth weight

Effect:
The player adds a D6 to their dice pool for a duration of 1D6 minutes. This applies to all attribute, skill, combat, magic, knowledge rolls, etc.
After the effect has expired, the player is easily irritable and prone to arguments.

    \paragraph{\trans{ingredients_label}}
    \begin{denseitemize}
        \item \textbf{Vial of flying lizard blood} (2 Vial), \textit{Components}.
        Mostly extracted from the veins of flying lizards domesticated by the O'Gru. It is no different from the blood of specimens living in the wild.

One vial contains 50 units of flying lizard blood.
        \item \textbf{Frost Lichen} (100 oth), \textit{Components}.
        A tough plant that thrives exclusively in the northern steppes and the coldest regions of the Tirakan Mountains.

It prefers to grow at the tree line and on barren, windswept rocky outcrops, where temperatures rarely rise above freezing even in summer. It is well known to dwarf prospectors from the north.

The lichen itself has no magic of its own, but it has an extreme cold-binding property. It absorbs the arctic cold of its surroundings and stores it for up to 3 days if the ambient temperature does not rise above 30°C. Within this period, the temperature of the plant remains that of its last location.

In alchemy, it therefore serves as a catalyst for stabilization, putting strong, volatile substances (such as blood or high-proof alcohol) into a state of “cold shock.”

Frost lichen appears as an inconspicuous, dense network in deep blue or white-gray. It lies like a crusted carpet on the stones and, at first glance, is hardly distinguishable from ice-covered rocks or frozen moss. It has no leaves and no flowers. 

It is often traded in the southern kingdoms of men, as it does not grow there but is essential for simple healing and strengthening potions.
        \item \textbf{Beet schnapps} (400 oth), \textit{Food / Provisions}.
        Beet schnapps is a typical product of human farmland in the more temperate zones of Tirakan, especially in the western and central kingdoms of the humans.

It is made from fermented and distilled sugar beets.

It is a cheap but high-proof spirit. In alchemy, it is used solely as a solvent and heat source to extract the aggressive or volatile properties of other substances (such as the bile of the flying lizard). Due to its purity and lack of complex ingredients, it is ideal as a simple base for mass carriers.

Beet schnapps is usually clear or slightly yellowish-cloudy and has a pungent smell of ethanol and a subtle, earthy sweetness. It burns when drunk and leaves a strong, unpleasant aftertaste.

Beet schnapps symbolizes the endurance and pragmatism of the people of Tirakan. While the elves have their “living water” and the dwarves their deep salt, humans rely on simple, readily available solutions.

It is a mass-produced item and an important commodity in the border regions and mercenary camps. A large part of the price of simple potions is accounted for by distillation and transport, not the ingredient itself.

For the high-ranking alchemists in the academies, beet schnapps is a sign of amateurism; they prefer more refined, less aggressive solvents. Village alchemists, on the other hand, use it because of its efficiency and availability.
    \end{denseitemize}


    \subsection*{Beturia's eternal rest}

    \begin{tabular}{@{}ll@{}}
        \textbf{\trans{difficulty_label}:} & medium \\
        \textbf{\trans{yield_label}:} & 1 \\
    \end{tabular}

    Preparation:
The brewing process begins by boiling water, valerian, and moonthorn berries. The mixture must first be boiled down to half its original volume. After removing it from the heat, the Nightslate dust is slowly stirred in until the potion turns a deep black color. The brew must then be left to mature in a sealed container for a full day (24 hours).

Preparation Time: 1 Stunde, 24 hours until ready

Quantity: 1 Application, 300 oth liquid, 320 oth weight

Effect:
After ingestion, the character immediately falls into a 24-hour state of suspended animation, during which their bodily functions are drastically reduced and they heal 2d6 wounds. To awaken from this near-comatose sleep, the character needs four successful rolls on Resistance in a row. Since shaking, loud noises, or attacks do not wake them, they can only be identified as alive by two successful rolls on Investigate; otherwise, they are considered dead.

    \paragraph{\trans{ingredients_label}}
    \begin{denseitemize}
        \item \textbf{Valeriana (Valeriana officinalis)} (4 Bunch), \textit{Components}.
        Valerian helps with insomnia and restlessness. Hops and lemon balm increase the effect of valerian and improve the taste.
        \item \textbf{Nightslate dust} (120 oth), \textit{Components}.
        Nightslate is no ordinary mineral, but the finely ground essence of ancient, jet-black rock veins.

This material is not mined, but must be extracted from the heart of mines that have not seen sunlight for eons, often accessible only through narrow tunnels in the coldest, most remote mountain ranges.

It is found exclusively in deep, abandoned mines or underground crypts, where the rock has been compacted over millennia by the constant pressure and absolute cold of the earth's interior. The veins shimmer slightly when caught in the glow of a torch, a sign of their almost unnatural purity.
        \item \textbf{Water} (500 oth), \textit{Components}.
        Cold, clear water.
    \end{denseitemize}


    \subsection*{Elixir of sweet slumber}

    \begin{tabular}{@{}ll@{}}
        \textbf{\trans{difficulty_label}:} & simple \\
        \textbf{\trans{yield_label}:} & 2 \\
    \end{tabular}

    Preparation:
First, heat the water in a suitable container and bring it to a boil. Once it is boiling vigorously, add the valerian and allow the mixture to boil for five minutes so that the active ingredients can be released. Then remove the container from the heat and allow the mixture to steep for another thirty minutes. Finally, strain the finished liquid and pour it into a clean container. Add honey for taste.

Quantity: 2 applications, 2x250 Oth liquid, 2x250 Oth weight

Effect:
After consumption, the character falls into a deep sleep within ten minutes, which lasts at least five hours. If the character is disturbed during this time or if someone tries to wake them up, they must make a resistance roll. If successful, the character wakes up immediately. If the roll fails, the character remains asleep, and the roll can be repeated the next time they are disturbed or someone tries to wake them again. However, it is important to note that if the character is shaken, attacked, or exposed to very loud noises, they will wake up immediately.

    \paragraph{\trans{ingredients_label}}
    \begin{denseitemize}
        \item \textbf{Valeriana (Valeriana officinalis)} (1 Bunch), \textit{Components}.
        Valerian helps with insomnia and restlessness. Hops and lemon balm increase the effect of valerian and improve the taste.
        \item \textbf{Water} (500 oth), \textit{Components}.
        Cold, clear water.
    \end{denseitemize}


    \subsection*{Elixir of elven power}

    \begin{tabular}{@{}ll@{}}
        \textbf{\trans{difficulty_label}:} & medium \\
        \textbf{\trans{yield_label}:} & 2 \\
    \end{tabular}

    Preparation:
The oil from the Schiller whale is slowly heated to body temperature together with the water in a glass bowl. It must never boil, otherwise the “shimmer” (the magical essence) will evaporate. The rock salt is sprinkled in until the liquid becomes milky.
The silver orchid leaves are then placed in the warm oil (maceration). They must rest there for 11 hours and slowly release their juices into the oil.
Finally, the elixir is filtered through a silk cloth. The remaining liquid should now have taken on the characteristic pearly glow of the whale.

Preparation time: 1 hour of active time (for precise temperature control) and 11 hours of resting time. 

Quantity: 2 applications, 2x25 Oth liquid, 2x25 Oth weight

Effect:
The player adds 2D6 to their dice pool for a duration of 2D6 minutes.
After this time, the character suffers the “shocked” 2 condition and is stunned for 2D6 minutes as their metabolism abruptly slows down after the highly potent oil wears off. Elven characters do not suffer this side effect.

    \paragraph{\trans{ingredients_label}}
    \begin{denseitemize}
        \item \textbf{Silberorchidee} (3 Petals), \textit{Components}.
        The silver orchid is considered the undisputed and deceptive “queen of southern flora.” It is a botanical marvel that is as beautiful as it is deadly to those who are fooled by its splendor.

Its leaves are not green, but have a dark gray, almost metallic color that shines like polished silver in the moonlight. The veins pulsate faintly in a pale violet when magic is nearby. The flower itself is large and cup-shaped, with snow-white petals that are razor-sharp at the edges. But the most disturbing thing about it is not its beauty, but its mobility: the plant stretches upward on exposed, muscular roots, which enable it to crawl slowly across the ground.

The silver orchid is found almost exclusively in the deep south of Tirakan, beyond the Iron Mountains. It can be found in the vast green steppes and along riverbanks, often camouflaged in the shade of the local flora. It often grows in disturbing proximity to the giant stone creatures that dwell in the passes.

When you approach the plant, it emits a glittering cloud of fine, silver dust. This is not harmless pollen, but a deadly attack. Anyone who inhales the dust is seized by severe coughing and shortness of breath. Within moments, black pockmarks form on the skin, and the victim falls into a deep, death-like unconsciousness. Once the victim is defenseless, the plant secretes a corrosive substance from its roots to slowly decompose and absorb its prey.

Despite these dangers, it is hunted because it is a powerful potentiator. In alchemy, the extracted nectar is used to increase the effects of other potions to the extreme. But processing it is risky: a mistake in distillation causes the magical energy to overload the body (similar to “silver death”), which is often fatal.

Ancient legends say that silver orchids came into being when the blood of a fallen star god dripped onto the earth in the First Age. The elves, on the other hand, call the flower “traitor's jewelry” and believe that it grows where reality has cracked and chaos seeps into the world.

“It looks like jewelry, crawls like a spider, and is worth more than my house. But be careful, boy: when you see the glitter, hold your breath and run. Before you realize you're getting sick, you'll already be its fertilizer.” – Marginal note in the records of the herb collector 'Three-Finger Hannes'
        \item \textbf{Schiller whale oil} (1 Vial), \textit{Components}.
        The blubber, the fatty tissue under the skin, is the reason why these majestic animals are hunted. It is no ordinary fat, but a storehouse of magical energy.

In its raw state, blubber is a tough, jelly-like mass that glows faintly blue. After refinement (melting and filtering), it becomes a clear, oily elixir that streaks like liquid mother-of-pearl.
Unlike the rancid blubber of ordinary whales, the blubber of the Schiller whale smells fresh, salty, and slightly metallic (like the air before a thunderstorm).

It is the best known means of binding volatile magic in potions (see Elixir of the Elven Watch).

When burned in lamps, it gives off a light that never produces soot and can make the invisible visible. Weapon oils made from this oil can injure spirits.

Since hunting them is extremely dangerous (iridescent whales rarely defend themselves, but are often protected by sea elementals or mermaids) and the animals are rare, the price is enormous.

The Ancatir consider hunting iridescent whales a sacrilege. They only use oil that comes from whales that have washed ashore naturally (“gift of the tides”). Alchemists who use “bloody oil” are often expelled from the city in elven enclaves.

The Schiller whale is one of the most fascinating and peace-loving giants of the seas around Tirakan. It is not just an animal, but a living anomaly closely connected to the magical currents of the oceans.

The iridescent whale resembles an earthly blue whale in shape, but is slimmer and has longer, almost wing-like side fins. What makes it special is its skin: it is not gray, but has a pearlescent, semi-transparent surface.
Depending on the incidence of light and the magical saturation of the environment, its skin refracts light into all colors of the spectrum. Hence the name. When a iridescent whale breaks the surface, it looks like a living rainbow rising out of the water.
Scholars believe that these whales not only feed on krill, but also absorb the light of the moon and stars when they come to the surface at night.

Iridescent whales avoid shallow coastal waters. They travel through the deep oceans, far away from the routes of merchant ships.

They are most often sighted in the Southern Ocean, in the cold currents far from the heat of the jungle, or in the mystical waters around abandoned island archives.

They travel in small family groups (pods). It is said that their song can calm storms or drive madness in those who listen to it for too long. 

We saw it at new moon. It glowed beneath the keel like a sunken city. When we threw the harpoons, the beast did not scream. It began to sing. A sound so deep that the wood of my ship splintered and two of my men simply jumped into the water, smiling. We killed it, yes. But the oil... it burns in the lamps of my cabin, and I swear I see the faces of those who jumped in the shadows.*
From the logbook of the whaler 'Haken-Ulf'
    \end{denseitemize}


    \subsection*{Lachsunder Seemagen}

    \begin{tabular}{@{}ll@{}}
        \textbf{\trans{difficulty_label}:} & simple \\
        \textbf{\trans{yield_label}:} & 3 \\
    \end{tabular}

    Preparation:
Take your small cauldron and fill it with clear spring water. Crush the dried mugwort and seaweed with a mortar and pestle until a coarse paste forms. Now add freshly grated ginger. Bring the mixture to a brief boil, then simmer over moderate heat until the liquid turns dark green. Strain the decoction through a linen cloth and pour it into a vial.

Preparation time:
20 minutes of active brewing time; no resting time required.

Quantity:
3 doses (liquid: 30 Oth, total weight 90 Oth).

Effect:
Instantly heals all ailments caused by seasickness (can also be taken preventively up to 2 hours before a sea voyage) and heals a wound after 2 hours.

    \paragraph{\trans{ingredients_label}}
    \begin{denseitemize}
        \item \textbf{Mugwort (Artemisia vulgaris)} (1 Bunch), \textit{Components}.
        A mugwort plant. The tops of the sprout are used to revive the digestion.
        \item \textbf{Seaweed from Lachsund  (Alga asgorania)} (10 oth), \textit{Components}.
        Eine zähe, dunkelgrüne bis bräunliche Meeresalge mit breiten, ledrigen Blättern, die an den Küsten von Lachsund im kalten Meerwasser gedeiht.

Sie schmeckt extrem salzig und verströmt einen intensiven Duft nach frischer Gischt und Jod.

In der alchemistischen Naturkunde dient er als Stabilisator für den Magen. Seine Inhaltsstoffe beruhigen den Gleichgewichtssinn und binden Giftstoffe im Verdauungstrakt. Für den alltäglichen Verzehr ist er nicht geeignet

Das Gewächs wird in großen Mengen an den felsigen Küsten von Asgoran angespült oder gezielt von den maritimen Sammlern der Siederei Falke aus dem Meer gefischt, getrocknet und verarbeitet. Auf den Fischmärkten von Lachsund ist er körbeweise erhältlich.
        \item \textbf{Lemon Ginger (Zingiber meridionale)} (5 oth), \textit{Components}.
        A tuberous, fibrous root of light brown color, whose juicy, yellow flesh releases a burning pungency and an intensely fresh citrus aroma when cut open.

In alchemy, it is considered a fiery vitalizer that stimulates blood circulation. It is used to stimulate the body’s vital fluids, alleviate stomach ailments, and drastically accelerate the absorption of healing elixirs into the bloodstream.

The plant thrives in warm, humid conditions and grows magnificently in the fertile, river-rich lands of O’grut, or is imported to Meridian from the southern oases via the humans’ extensive trade networks.
        \item \textbf{Water} (90 oth), \textit{Components}.
        Cold, clear water.
    \end{denseitemize}


    \subsection*{Panthomeic Pick-Me-Up}

    \begin{tabular}{@{}ll@{}}
        \textbf{\trans{difficulty_label}:} & medium \\
        \textbf{\trans{yield_label}:} & 2 \\
    \end{tabular}

    Preparation:
First, you must grind the hard Panthomean pine cones into a fine powder using a mortar. Bring this powder to a boil in clear water over a blue, extremely hot flame until the water turns emerald green. Now you must add fresh bloodwort to neutralize the bitter compounds in the cones. The chamomile, lightly crushed, gives the potion its energizing effect.

Preparation time:
2 hours of intense boiling, followed by half a night of steeping in a cool cellar.

Quantity:
2 doses. 200 Oth of liquid, total weight 200 Oth.

Effect:
Temporarily eliminates sleep deprivation. The user ignores all penalties due to fatigue and also receives a substantial bonus of +2 dice to Perception and Alertness for 1W6 hours. However, as soon as the effect wears off, the body inevitably falls into a deep, dreamless sleep.

    \paragraph{\trans{ingredients_label}}
    \begin{denseitemize}
        \item \textbf{Panthomic pine cones} (2 Piece), \textit{Components}.
        Panthomean pine cones are the woody, resin-rich fruit clusters of the rare Panthomean mountain pine.

These cones are covered by a hard, almost metallic-looking layer of scales that encloses a deep-purple, highly sticky balsam inside. They possess no inherent magical energy, but are known in alchemy for their powerful preservative, structure-strengthening, and heat-generating properties.

These cones grow exclusively in the icy, storm-lashed heights of the Panthomean peaks, far above the tree line in the rugged highlands.

Harvesting them is arduous, as the cones can only be harvested by hunters or the Hil-JabalJarh.

Occasionally, dwarven mountain caravans bring isolated specimens to markets throughout Tirakan.
        \item \textbf{Common Bloodweed} (20 oth), \textit{Components}.
        Bloodweed is a ground-level plant that is particularly striking due to its fleshy, deep red leaves.

The fine veins on the leaf surface glow in a rich scarlet red, almost as if real blood were pulsing through them. When a leaf is crushed, a sticky, sweet-smelling sap emerges that stains the fingers for days. It is not a magical plant in the classic sense. It draws its power from the iron-rich soil and the pale light of the dense forests.

You will usually find bloodweed in shady, damp places. It prefers the foot of old oak trees or the immediate vicinity of rotting undergrowth in deep forests. An inattentive traveler often mistakes it for common purple sorrel, but a trained alchemist will notice the small, pearl-like dewdrops that always collect at the edges of the leaves.

It is best harvested in the early morning hours, before the sun breaks through the canopy. Only the outer leaves are cut to preserve the root.
        \item \textbf{Camomile (Matricaria recutita)} (1 Bunch), \textit{Components}.
        Chamomile is one of the oldest medicinal plants and was already used in the Middle Ages. The flowers have a healing and soothing effect. Externally, chamomile can be used for inflammation of the gums, skin or mucous membrane. Taken internally, it is effective for gastrointestinal disorders. Rinsing and inhalation are also widely used.
        \item \textbf{Water} (130 oth), \textit{Components}.
        Cold, clear water.
    \end{denseitemize}


    \subsection*{Snakeoil}

    \begin{tabular}{@{}ll@{}}
        \textbf{\trans{difficulty_label}:} & simple \\
        \textbf{\trans{yield_label}:} & 1 \\
    \end{tabular}

    Preparation:
First, melt the fat (any kind, as long as it's fatty!) in a small iron pot over moderate heat. Make sure it doesn't burn. Once the fat is liquid, stir in the mint decoction. This serves solely to mask the pungent smell of the fat and give the tincture a “medicinal” green color. A dash of beet schnapps ensures that the mixture burns in your throat. But what burns also helps. Finally, shake the emulsion vigorously until it looks cloudy and mysterious.

Preparation time: 30 minutes of active stirring, 1 hour to cool and settle.

Quantity: 1 application (liquid: 5 Oth / weight: 6 Oth including bottle).

Effect: Immediately heals 1 wound when applied. The patient feels briefly invigorated but suffers from the terrible aftertaste.

    \paragraph{\trans{ingredients_label}}
    \begin{denseitemize}
        \item \textbf{Lard} (10 oth), \textit{Components}.
        It is fat. Of plant or animal origin. It is used for frying, refining food, or lubricating.
        \item \textbf{Mint} (1 Bunch), \textit{Components}.
        Mint is widespread in Tirakan, but among those who study herbal medicine, it is considered an indispensable staple for the mind and body.

Occurrence \& growth:
Mint is usually found in damp, semi-shaded locations. It grows rampantly on the banks of streams, in enchanted forest clearings, or in the herb gardens of wise healers. 

It is considered to be soothing for the stomach, has a cooling effect, and relieves sore throats. It is used in ointments, tinctures, and teas. It is also often used to soften the strong taste of game meat or to give cheap thin beer a fresh note.
        \item \textbf{Water} (10 oth), \textit{Components}.
        Cold, clear water.
    \end{denseitemize}


    \subsection*{Brew of Uncontrolled Zeal (Goblin Variant)}

    \begin{tabular}{@{}ll@{}}
        \textbf{\trans{difficulty_label}:} & simple \\
        \textbf{\trans{yield_label}:} & 1 \\
    \end{tabular}

    Preparation:
Roughly tear the dried leaves of the guild herb with your fingers and toss them into an uncleaned iron cauldron along with cheap beet schnapps. While clapping loudly and rhythmically, bring the mixture to a brief, vigorous boil until a pungent, nutty steam rises. Immediately remove from the heat and pour the warm liquid into a leather wineskin.

Preparation time:
15 minutes of active work (no resting time).

Quantity: 1 dose (10 Oth of liquid, weight 12 Oth).

Effect: This brew affects only members of the goblin races. If another creature drinks this potion, it feels slightly more energetic and is full of drive.

The goblin, however, falls into a work frenzy for 1d6×10 minutes. It gains +2 dice on all craft or physical skill checks (e.g., Mechanics, Mining, Climbing). However, all dexterity or logic checks (Logic, Deception, and Cheating) are made +3 more difficult during this time, as their mind jumps back and forth completely uncoordinated. They cannot control themselves and can hardly be stopped, except by strong arms or restraints.

    \paragraph{\trans{ingredients_label}}
    \begin{denseitemize}
        \item \textbf{Gildenkraut (Vigilia mercatoris)} (5 oth), \textit{Components}.
        Gildherb is a tough, almost inconspicuous plant that only catches the trained eye of an expert.

It is characterized by jagged, leathery leaves of a deep green color, whose thick veins have a slight reddish sheen. Its flowers are tiny, pale yellow, and close immediately once the sun passes its zenith.

The herb is entirely mundane and possesses not the slightest spark of the arcane. Its value lies rather in its extremely high, natural concentration of alkaloids, particularly pure caffeine. When the tough leaves are crushed between the fingers, they give off a sharp, earthy scent reminiscent of heavily roasted nuts.

Location and Habitat:
This plant is relatively rare, as it requires very specific, adverse conditions to produce its stimulating sap. It thrives exclusively in barren, mineral-rich soils and avoids the damp moss of deep forests.

Gildherb is found primarily in the rocky crevices of ancient trade routes, in the crumbling stone mortar of abandoned guild outposts, or on the steep, sun-baked slopes of the Toran Highlands. It needs hard rock and plenty of wind to prevent it from withering away.
        \item \textbf{Beet schnapps} (10 oth), \textit{Food / Provisions}.
        Beet schnapps is a typical product of human farmland in the more temperate zones of Tirakan, especially in the western and central kingdoms of the humans.

It is made from fermented and distilled sugar beets.

It is a cheap but high-proof spirit. In alchemy, it is used solely as a solvent and heat source to extract the aggressive or volatile properties of other substances (such as the bile of the flying lizard). Due to its purity and lack of complex ingredients, it is ideal as a simple base for mass carriers.

Beet schnapps is usually clear or slightly yellowish-cloudy and has a pungent smell of ethanol and a subtle, earthy sweetness. It burns when drunk and leaves a strong, unpleasant aftertaste.

Beet schnapps symbolizes the endurance and pragmatism of the people of Tirakan. While the elves have their “living water” and the dwarves their deep salt, humans rely on simple, readily available solutions.

It is a mass-produced item and an important commodity in the border regions and mercenary camps. A large part of the price of simple potions is accounted for by distillation and transport, not the ingredient itself.

For the high-ranking alchemists in the academies, beet schnapps is a sign of amateurism; they prefer more refined, less aggressive solvents. Village alchemists, on the other hand, use it because of its efficiency and availability.
    \end{denseitemize}


    \subsection*{Falk's Tran Ointment}

    \begin{tabular}{@{}ll@{}}
        \textbf{\trans{difficulty_label}:} & medium \\
        \textbf{\trans{yield_label}:} & 2 \\
    \end{tabular}

    Preparation:
Take a small iron cauldron and heat the Schillerwal oil over moderate heat until it becomes completely liquid and clear.

While stirring constantly and rhythmically, add the dried and finely ground ribwort plantain to give the ointment its wound-healing and skin-tightening properties. As soon as the mixture takes on a slight bluish sheen, remove the cauldron from the coals and stir in the purified fat while slowly pouring the mixture into earthenware pots, where it will solidify into a thick paste.

Preparation time:
1 hour of active simmering and stirring, 3 hours of resting time to allow it to solidify.

Quantity:
2 applications, 70 Oth of pure ointment mixture.

Effect:
Application takes 1 minute. Grants the user 1 point of wound healing (heals any wound) for the next 2 hours, as well as the protection units N (Normal Protection) and B (Bleeding Protection) for the next battle. Additionally, the user is immune to the effects of natural cold during this time.

    \paragraph{\trans{ingredients_label}}
    \begin{denseitemize}
        \item \textbf{Schiller whale oil} (50 oth), \textit{Components}.
        The blubber, the fatty tissue under the skin, is the reason why these majestic animals are hunted. It is no ordinary fat, but a storehouse of magical energy.

In its raw state, blubber is a tough, jelly-like mass that glows faintly blue. After refinement (melting and filtering), it becomes a clear, oily elixir that streaks like liquid mother-of-pearl.
Unlike the rancid blubber of ordinary whales, the blubber of the Schiller whale smells fresh, salty, and slightly metallic (like the air before a thunderstorm).

It is the best known means of binding volatile magic in potions (see Elixir of the Elven Watch).

When burned in lamps, it gives off a light that never produces soot and can make the invisible visible. Weapon oils made from this oil can injure spirits.

Since hunting them is extremely dangerous (iridescent whales rarely defend themselves, but are often protected by sea elementals or mermaids) and the animals are rare, the price is enormous.

The Ancatir consider hunting iridescent whales a sacrilege. They only use oil that comes from whales that have washed ashore naturally (“gift of the tides”). Alchemists who use “bloody oil” are often expelled from the city in elven enclaves.

The Schiller whale is one of the most fascinating and peace-loving giants of the seas around Tirakan. It is not just an animal, but a living anomaly closely connected to the magical currents of the oceans.

The iridescent whale resembles an earthly blue whale in shape, but is slimmer and has longer, almost wing-like side fins. What makes it special is its skin: it is not gray, but has a pearlescent, semi-transparent surface.
Depending on the incidence of light and the magical saturation of the environment, its skin refracts light into all colors of the spectrum. Hence the name. When a iridescent whale breaks the surface, it looks like a living rainbow rising out of the water.
Scholars believe that these whales not only feed on krill, but also absorb the light of the moon and stars when they come to the surface at night.

Iridescent whales avoid shallow coastal waters. They travel through the deep oceans, far away from the routes of merchant ships.

They are most often sighted in the Southern Ocean, in the cold currents far from the heat of the jungle, or in the mystical waters around abandoned island archives.

They travel in small family groups (pods). It is said that their song can calm storms or drive madness in those who listen to it for too long. 

We saw it at new moon. It glowed beneath the keel like a sunken city. When we threw the harpoons, the beast did not scream. It began to sing. A sound so deep that the wood of my ship splintered and two of my men simply jumped into the water, smiling. We killed it, yes. But the oil... it burns in the lamps of my cabin, and I swear I see the faces of those who jumped in the shadows.*
From the logbook of the whaler 'Haken-Ulf'
        \item \textbf{Ribwort plantain (Plantago lanceolata)} (1 Bunch), \textit{Components}.
        The pointed, narrow leaves of ribwort plantain are used as a syrup or also as a tea for colds. Ribwort can also be crushed and ground and applied to wounds or insect bites, where it has a cooling effect. The plant is also used for diarrhea.
        \item \textbf{Lard} (20 oth), \textit{Components}.
        It is fat. Of plant or animal origin. It is used for frying, refining food, or lubricating.
    \end{denseitemize}


    \subsection*{Potion of Protection}

    \begin{tabular}{@{}ll@{}}
        \textbf{\trans{difficulty_label}:} & simple \\
        \textbf{\trans{yield_label}:} & 1 \\
    \end{tabular}

    Preparation:
First, bring the spring water to a boil in a copper kettle until the steam fogs up the windows. Add the dried mugwort leaves. As soon as the water turns a deep green color, stir in the pine resin. Be careful here: stir slowly and always in a clockwise direction, otherwise the resin will clump and you will end up with sticky slime instead of liquid protection! Finally, finely crush the elecampane root and add it to bind the vitality. Filter the mixture through a fine linen cloth into a vial. 

Preparation time: 45 minutes of active brewing time, followed by 1 hours of resting time to cool and bind the essences.

Quantity:
1 application (approx. 8 Oth liquid, weight of vial including contents: 10 Oth).

Effect:
When applied, the character receives a 1D3 boost to their armor. Each of these boosts can be used as Normal Protection (N). This means that a boost can be used to completely prevent a normal hit.

    \paragraph{\trans{ingredients_label}}
    \begin{denseitemize}
        \item \textbf{Water} (50 oth), \textit{Components}.
        Cold, clear water.
        \item \textbf{Mugwort (Artemisia vulgaris)} (1 Bunch), \textit{Components}.
        A mugwort plant. The tops of the sprout are used to revive the digestion.
        \item \textbf{Inula (Inula helenium)} (1 Bunch), \textit{Components}.
        This medicinal plant from the Middle Ages is no longer widely used in modern times. Its application improves digestion, and it is believed to have a preventive effect against colon cancer.
        \item \textbf{Pine Resin} (10 oth), \textit{Components}.
        The protective blood of the tree. When the rough bark of a pine tree is damaged, whether by the bite of a wild animal or the clumsy axe of a lumberjack, this tough, golden-yellow liquid oozes out.

It flows slowly, almost sluggishly, filling the air with its distinctive, tart scent.

It is used to seal brittle corks, give torches a strong flame, or thicken simple wound ointments.

Once it dries in the air, it becomes rock hard and closes the tree's wound like a natural plug. 

Once you get this stuff on your hands, the only way to get rid of the glue is with a lot of grease or a lot of patience.
    \end{denseitemize}


    \subsection*{Elixir of Chronos}

    \begin{tabular}{@{}ll@{}}
        \textbf{\trans{difficulty_label}:} & difficult \\
        \textbf{\trans{yield_label}:} & 1 \\
    \end{tabular}

    Preparation:
Bring the pure spring water to a boil in a heavy iron cauldron. Add the obsidian dust and stir continuously until the water takes on a cloudy, gray color. Now add the dried guild herb and the valerian. The dust binds the plant fibers and slowly settles to the bottom during the three-hour, vigorous boil, releasing the concentrated alkaloids into the liquid.

The mixture must be filtered through a triple layer of linen cloth. If obsidian dust remains in the potion, it weighs heavily on the consumer’s gut, causing 1D6 hours of exhaustion. All checks are made with a +1 penalty during this time.

Preparation time:
3 hours of active, highly concentrated work. 12 hours of rest in absolute darkness and silence.

Quantity: 1 dose (approx. 15 Oth of pure essence, total weight 15 Oth).

Effect:
For the next 6 hours, the consumer enters a state of absolute, supernatural hyper-focus.
The character gains a permanent bonus of +2 dice on all checks requiring Logic, Diligence, Perception, or Knowledge.

Additionally, for the duration of the effect, the character becomes completely immune to the Dazed condition as well as to all attempts at mental influence or sleep spells (since the brain is simply working too fast).

As soon as the 6 hours elapse, the nervous system collapses abruptly. The character instantly falls into a comatose deep sleep for 1D6+2 hours, from which they cannot be awakened under any circumstances (not even damage will wake them). They awaken afterward somewhat confused and cannot remember what happened during the 6-hour duration.

However, the deep sleep allows the character to make a Rest roll.

    \paragraph{\trans{ingredients_label}}
    \begin{denseitemize}
        \item \textbf{Gildenkraut (Vigilia mercatoris)} (30 oth), \textit{Components}.
        Gildherb is a tough, almost inconspicuous plant that only catches the trained eye of an expert.

It is characterized by jagged, leathery leaves of a deep green color, whose thick veins have a slight reddish sheen. Its flowers are tiny, pale yellow, and close immediately once the sun passes its zenith.

The herb is entirely mundane and possesses not the slightest spark of the arcane. Its value lies rather in its extremely high, natural concentration of alkaloids, particularly pure caffeine. When the tough leaves are crushed between the fingers, they give off a sharp, earthy scent reminiscent of heavily roasted nuts.

Location and Habitat:
This plant is relatively rare, as it requires very specific, adverse conditions to produce its stimulating sap. It thrives exclusively in barren, mineral-rich soils and avoids the damp moss of deep forests.

Gildherb is found primarily in the rocky crevices of ancient trade routes, in the crumbling stone mortar of abandoned guild outposts, or on the steep, sun-baked slopes of the Toran Highlands. It needs hard rock and plenty of wind to prevent it from withering away.
        \item \textbf{Valeriana (Valeriana officinalis)} (1 Bunch), \textit{Components}.
        Valerian helps with insomnia and restlessness. Hops and lemon balm increase the effect of valerian and improve the taste.
        \item \textbf{Obsidian (Dust)} (5 oth), \textit{Components}.
        Obsidianstaub ist ein mundaner, mineralischer alchemistischer Grundstoff, der aus vulkanischem Glas gewonnen wird.

Er besitzt keinerlei innewohnende magische Kraft, ist jedoch aufgrund seiner extremen Härte, seiner scharfen Kanten im mikroskopischen Bereich und seiner enormen Hitzebeständigkeit ein unschätzbarer mechanischer Katalysator.

In der Alchemie wird er genutzt, um Unreinheiten und Bitterstoffe aus kochenden Suden zu binden und feine Pflanzenessenzen zu zwingen, sich abzusetzten.

Obsidian entsteht dort, wo feurige Magma rasch an der frischen Luft oder durch Wasser abkühlt. Man findet das vulkanische Glas in den tiefen, erstarrten Lavaseen der Unterwelt oder an den rauen, vulkanischen Bruchstellen des Hochlands.
Er kann in schweren Mörsern zu einem mehlfeinen, grauen Staub zerstoßen werden.

Die unangefochtenen Hauptlieferanten für den besten Obsidianstaub in Tirakan sind die Zwerge der Xordai. Insbesondere die Kolonie Kor-Vamak, die direkt in einem erstarrten Lavasee errichtet wurde, hat sich auf den Abbau und die Verarbeitung dieses Materials spezialisiert.

Das dort ansässige Syndikat exportiert den Staub über ihre gut bewachten Handelskarawanen bis an die Oberfläche.
    \end{denseitemize}


    \subsection*{Felsenfeuer Distillate}

    \begin{tabular}{@{}ll@{}}
        \textbf{\trans{difficulty_label}:} & simple \\
        \textbf{\trans{yield_label}:} & 1 \\
    \end{tabular}

    Preparation:
The preparation requires meridian-like
precision. First, the marshmallow root is steeped in cold water
and left to steep for exactly two hours before the liquid
is strained and heated. In the meantime, comfrey and chamomile are ground into a fine powder in a mortar. The heated marshmallow liquid is mixed with the powder and distilled several times until a clear essence is obtained.

Preparation time:
2 hours of steeping time for the cold
extraction of the marshmallow root, followed by 1 hour of active time
at the distillation apparatus.

Quantity:
1 application. Approx. 200 Oth of liquid, total
weight 800 Oth (including ceramic bottle).

Effect:
The concentrated active ingredients accelerate cell regeneration and heal 1W3 wounds.

    \paragraph{\trans{ingredients_label}}
    \begin{denseitemize}
        \item \textbf{Water} (200 oth), \textit{Components}.
        Cold, clear water.
        \item \textbf{Camomile (Matricaria recutita)} (50), \textit{Components}.
        Chamomile is one of the oldest medicinal plants and was already used in the Middle Ages. The flowers have a healing and soothing effect. Externally, chamomile can be used for inflammation of the gums, skin or mucous membrane. Taken internally, it is effective for gastrointestinal disorders. Rinsing and inhalation are also widely used.
        \item \textbf{Comfrey (Symphytum officinale)} (50), \textit{Components}.
        Comfrey stimulates blood circulation, bruises, hematomas and sprains disappear faster. Comfrey accelerates the regeneration of cells.
        \item \textbf{Marshmallow (Althaea officinalis)} (100), \textit{Components}.
        The root of this medicinal plant is used. This is prepared cold and must infuse for about two hours. Only after infusion, the liquid is strained and then heated. The substances provide protection for the mucous membranes and have an anti-irritant effect. A helpful medicinal plant for gastrointestinal problems and a cough.
    \end{denseitemize}


    \subsection*{Mist of Thenon}

    \begin{tabular}{@{}ll@{}}
        \textbf{\trans{difficulty_label}:} & medium \\
        \textbf{\trans{yield_label}:} & 1 \\
    \end{tabular}

    Preparation:
Take a clean glass vial and drip the secretion from the Kinstarchel bone extract onto the bottom to ensure an explosive release upon impact.

Meanwhile, boil the dried ribwort plantain leaves in pure spring water until the liquid takes on a creamy, slimy consistency. While stirring constantly in a counterclockwise direction, add the ground amber. Add the resin; it will bind the vapors. Immediately seal the vial with liquid wax while the vapor is still hot.

Preparation time:
3 hours of active alchemy, followed by 1 hour of cooling in the dark.

Quantity:
1 use (vial holds approx. 50 Oth of liquid)

Effect:
When broken, instantly creates a stationary cloud of mist within a 3-meter radius that lasts for 5 rounds. All creatures within it are afflicted with the Blind condition, and all vision is severely restricted (attacks and perception checks from outside to inside and vice versa are considered difficult checks with a +2 DC).

    \paragraph{\trans{ingredients_label}}
    \begin{denseitemize}
        \item \textbf{Kinstarchel Secretion} (10 oth), \textit{Components}.
        A highly reactive liquid produced from the bones of the adorable Kinstarchel.
        \item \textbf{Ribwort plantain (Plantago lanceolata)} (1 Bunch), \textit{Components}.
        The pointed, narrow leaves of ribwort plantain are used as a syrup or also as a tea for colds. Ribwort can also be crushed and ground and applied to wounds or insect bites, where it has a cooling effect. The plant is also used for diarrhea.
        \item \textbf{Amber} (20 oth), \textit{Components}.
        A smooth, oval-shaped amber with a warm golden hue. Its polished surface is slightly transparent and reflects light in a fascinating way. The hand-sized stone looks like a natural talisman due to its curved shape.
        \item \textbf{Pine Resin} (10 oth), \textit{Components}.
        The protective blood of the tree. When the rough bark of a pine tree is damaged, whether by the bite of a wild animal or the clumsy axe of a lumberjack, this tough, golden-yellow liquid oozes out.

It flows slowly, almost sluggishly, filling the air with its distinctive, tart scent.

It is used to seal brittle corks, give torches a strong flame, or thicken simple wound ointments.

Once it dries in the air, it becomes rock hard and closes the tree's wound like a natural plug. 

Once you get this stuff on your hands, the only way to get rid of the glue is with a lot of grease or a lot of patience.
    \end{denseitemize}


    \subsection*{Potion of Might}

    \begin{tabular}{@{}ll@{}}
        \textbf{\trans{difficulty_label}:} & difficult \\
        \textbf{\trans{yield_label}:} & 1 \\
    \end{tabular}

    Preparation: First, the bile of a mountain wyvern must be slowly heated in a crucible made of hardened steel until the corrosive vapors have evaporated and a golden-yellow base substance remains. The grated claw pieces of the rock troll are added to this base, which bind the raw physical power. While stirring constantly, the dried Tirakan moss must now be added, which serves as a catalytic conductor. As soon as the mixture begins to pulsate slightly, the silver shavings are sprinkled in. The potion must now simmer over a fire for exactly one hour until it becomes thick and takes on a metallic sheen. After resting for one hour, it is ready for use.

Preparation time: 3 hours (active brewing), 12 hour resting time

Quantity: 1 application, 200 Oth liquid, 250 Oth weight

Effect: After ingestion, the character feels their muscles harden and their senses sharpen. The number of dice rolled by the player is doubled. The potion lasts for 2D6 minutes. The user also ignores all penalties due to exhaustion or wounds and receives temporary armor protection of 2 against normal damage (2xR) physical damage (this is added to existing armor).

As soon as the effect wears off, the magic takes its toll: the character immediately suffers 2D6+6 rounds of exhaustion and is “dazed” during this period (+2 penalty on all minimum rolls).

    \paragraph{\trans{ingredients_label}}
    \begin{denseitemize}
        \item \textbf{Mountain Wyvern Bile} (200 oth), \textit{Components}.
        Bile is a highly viscous, bright golden yellow to poison green liquid stored in the gallbladder of mountain wyverns.

It has a pungent, sulfurous odor with a hint of burnt copper. Even inhaling the pure vapors can burn the nasal mucous membranes.

In Tiraka, it is believed that bile contains the “essence of unquenched hunger.” In alchemy, it is used as a catalyst to forcibly fuse other ingredients that would normally repel each other.

Extracting bile is a difficult undertaking for an alchemist, as it requires the utmost precision under adverse conditions.

The bile must be extracted within 1D6 hours after the creature's death. After that, the gallbladder begins to decompose and the liquid loses its alchemical potency.

Surgical instruments made of hardened steel or special ceramic knives are required. Simple iron would be corroded by the acid within seconds.

The procedure is as follows: The carcass must be secured on its back.
A deep cut below the sternum exposes the liver.
The gallbladder is a bulging, pulsating sac. It must be clamped at the top before being carefully cut out. A successful Dexterity (or Medicine) roll against DC 8 is required.

If the sac bursts, the harvester immediately suffers 1D6 damage from chemical burns, and the ingredient is irretrievably lost. 

Wyvern bile cannot be stored in normal glass vials, as it will eventually “blind” the glass and cause it to become brittle. Experienced adventurers use lead-lined clay jugs or pure quartz vessels to safely transport the substance home.
        \item \textbf{Rock Moss} (2 Bunch), \textit{Components}.
        Tirakan moss, often called “shadow flora” or “shadow velvet” by mountain peoples, usually grows near mountain wyvern colonies. It can be found in rock crevices and small caves. It is said that it may purify air in narrow caves.

The moss grows in dense, sponge-like cushions. Its color is a deep, almost unnatural dark purple that glows in a soft, pulsating indigo when touched or exposed to air currents (bioluminescence).
When touched, it leaves a slightly sticky, metallic-smelling film on the skin.

Its main function in potions is grounding. It prevents energies or raw forces from dissipating.

When chewed raw, it has a strong pain-relieving and fever-reducing effect, but an overdose can lead to a dangerous slowing of the heartbeat.

The elders claim that the moss absorbs the whispers of the mountains. If you press your ear against a moss cushion long enough, you can hear the voices of your ancestors or the mountain growling with hunger.
        \item \textbf{Rock troll claws (Lithocrinus tirakanis)} (40 oth), \textit{Components}.
        The claws of a rock troll are not claws in the biological sense, but mineralized growths made of hardened keratin and concentrated ores.

They reflect the unbridled physical regenerative power of the trolls, whose skin has entered into a symbiosis with the rock of the Shadow Rocks.

A single, intact claw is about the size of a short sword.
Their color ranges from deep gray to obsidian black, often with layers reminiscent of slate.

They are so hard that they spark when they hit metal. Ordinary blades usually dull immediately when used on them.

In alchemy, the entire claw is almost never used, but rather a processed form. The claw must be laboriously worked with diamond files or hardened steel chisels. 

Finely grated shavings or dust are used as ingredients. This dust is heavy, ash-colored, and glitters when exposed to light.
The dust must be extremely fine. Chips that are too coarse will not dissolve in the potion and can seriously injure the user's esophagus when consumed. Alternatively, it is advisable to strain the liquid through a sieve after brewing.
        \item \textbf{silver shavings} (10 oth), \textit{Components}.
        In Tiraka, silver shavings are usually obtained as a by-product in forges or during the manufacture of jewelry.

For alchemical purposes, they are often purified in fire to stabilize the magical currents in potions as purified silver.
Silver shavings serve as an energetic anchor. They prevent the unstable components from “tearing apart” the potion during the brewing process.
    \end{denseitemize}


    \subsection*{Meridian Clarity}

    \begin{tabular}{@{}ll@{}}
        \textbf{\trans{difficulty_label}:} & medium \\
        \textbf{\trans{yield_label}:} & 2 \\
    \end{tabular}

    Preparation:
First, the silver thistle must be boiled in pure water and strained through the finest charcoal powder to activate the active ingredients and retain only those that are needed.

Next, the moonberries are carefully stirred in. Caution: The moonberries must not be crushed; they should simply steep for 12 hours.

Preparation time:
4 hours of active distillation and filtration using the apparatus, after which the brew must rest for exactly 12 hours.

Quantity:
2 applications. 400 Oth of liquid, total weight 600 Oth (including thick-walled glass vial).

Effect:
Neutralizes toxins in the bloodstream and sharpens the senses to protect the mind from hallucinations. Consuming one dose instantly neutralizes the “Poisoned” status. Additionally, the user gains a bonus of +1 die to any resistance checks against mind-altering effects or hallucinations for the next 1W6 hours.

    \paragraph{\trans{ingredients_label}}
    \begin{denseitemize}
        \item \textbf{Charcoal} (100 oth), \textit{Components}.
        A substance of pure, deep black color, obtained through the controlled carbonization of hard wood in the absence of air. Charcoal.

In alchemy, it serves as a universal filtering medium. Whether incorporated into recipes or ground into the finest powder, it absorbs impurities, cloudiness, and unwanted bitter substances from boiling liquids like a sponge.

It is widely available throughout Tirakan from charcoal burners, blacksmiths, and ordinary shopkeepers in large quantities for just a few deut or guilders at the city markets.
        \item \textbf{Silver Thistle (Carlina acaulis)} (2 Bunch), \textit{Components}.
        A hardy plant that grows close to the ground, with a silvery, shiny, spiky flower head that did not wither even in freezing temperatures.

In alchemy, it is considered a purifying catalyst whose essence binds toxins in the blood, inhibits inflammation, and strengthens the body’s natural defenses.

It thrives primarily in the barren, sun-drenched mountain meadows and rocky slopes of the Toran Highlands. However, it can also be found in more temperate regions throughout Tirakan.
        \item \textbf{Moonthorn Berry} (50 oth), \textit{Components}.
        The moonthorn berry is a gift from the deepest night. It grows as a ground-covering shrub whose delicate tendrils and deep green leaves are protected by striking, short thorns.

It is found exclusively in places that rarely see the light of the sun, usually deep in ancient forests or near damp grotto and cave entrances. It only reveals its true splendor under the light of the full moon, when its small, berry-like fruits glow in a mysterious, dull blue, almost as if they had swallowed the light of the celestial sphere itself.

The berry is notorious for its strong sedative effect. In small doses, it has a calming effect, but when concentrated in a potion, its essence can numb the mind and put the body into a state of deep, dreamless stillness. Gathering them is risky, as their thorns can cause temporary itching when touched.
        \item \textbf{Water} (350 oth), \textit{Components}.
        Cold, clear water.
    \end{denseitemize}


    \section{Poisons and Drugs}



    \subsection*{Linya's Deep Sleep}

    \begin{tabular}{@{}ll@{}}
        \textbf{\trans{difficulty_label}:} & difficult \\
        \textbf{\trans{yield_label}:} & 1 \\
    \end{tabular}

    Preparation:
Grind the dried black poppy seeds together with the valerian roots in a stone mortar into a uniform, black powder. Dissolve this in 150 Oth of pure spring water and let it steep undisturbed for two hours.

Strain the plant residue through a fine linen cloth into the small iron kettle. Add the lemon balm leaves to stabilize the organic ether.

Heat the kettle (preferably on the brewer’s table) while stirring constantly until the vapors take on a bluish sheen. Collect the condensate in a vial.

Preparation time: 2 hours of active resting time (steeping), 3 hours of active boiling and distillation.

Quantity: 1 dose (approx. 50 Oth of liquid)

Effect: If the essence is ingested or injected into the bloodstream via a weapon with a poison channel, the victim instantly falls into a state of unconsciousness.

The victim must roll for Willpower at the start of each of its combat rounds. Only when the number of successes (rolls of 5+; exploding 6s allowed) reaches or exceeds exactly 5 does the creature awaken. External stimuli such as shaking or noise have no effect; only actual damage suffered grants an immediate bonus of +2 dice to the saving throw.

    \paragraph{\trans{ingredients_label}}
    \begin{denseitemize}
        \item \textbf{Death Poppy} (2 Bunch), \textit{Components}.
        The poppy of the dead is cultivated exclusively on the asgoran island of Linya, the plant can be cultivated nowhere else. The plant is a poppy-like flower about twenty fingers high, which is partially colored black. 

The poppy of the dead is used for all kinds of rituals and potions, which makes the plant a real export of Asgoran.
        \item \textbf{Valeriana (Valeriana officinalis)} (1 Bunch), \textit{Components}.
        Valerian helps with insomnia and restlessness. Hops and lemon balm increase the effect of valerian and improve the taste.
        \item \textbf{Lemon balm (Melissa officinalis)} (1 Bunch), \textit{Components}.
        Lemon balm has always been used as a medicinal herb in medicine. It is effective against headaches, nervousness, insomnia and gastrointestinal complaints. In addition, an infusion with lemon balm brings relaxation.
        \item \textbf{Water} (150 oth), \textit{Components}.
        Cold, clear water.
    \end{denseitemize}


    \subsection*{Morphine}

    \begin{tabular}{@{}ll@{}}
        \textbf{\trans{difficulty_label}:} & difficult \\
        \textbf{\trans{yield_label}:} & 1 \\
    \end{tabular}

    Preparation:
As a natural painkiller, morphine is primarily found in opium poppies (Papaver somniferum). It is the most important component of opium, which is obtained by drying the milky sap from the unripe seed capsules. To isolate the active ingredient, the aqueous extract is treated with calcium chloride. This separates the calcium meconic acid. The remaining liquid is then evaporated to obtain morphine and codeine in the form of hydrochloride salts.

Preparation time: 60 minutes for extraction, followed by 1 hour to separate the morphine.

Quantity: 1 application, approx. 50 ml liquid, weight: 50 g.

Effect:
Morphine acts directly on the brain, where it alleviates pain perception and the transmission of pain signals. It also suppresses the urge to cough and has a calming effect on the body, which can also lower blood pressure and pulse rate.

For people who are not used to the drug, even 100 mg can lead to severe poisoning. However, those who take the drug on a long-term basis can often tolerate significantly higher doses. Morphine also has a high potential for addiction. A common misconception is that morphine hastens death. There is no scientific evidence to support this. Often, side effects such as severe drowsiness or confusion are mistakenly interpreted as signs of imminent death.

    \paragraph{\trans{ingredients_label}}
    \begin{denseitemize}
        \item \textbf{Calcium Chloride} (50 ml (milliliter)), \textit{Components}.
        Calcium chloride is a compound of calcium and chlorine with the formula CaCl\_2.

Chemically speaking, the calcium in it has an oxidation state of +2 and the chlorine has an oxidation state of -1. This substance is found in nature mainly dissolved in salt brines. When calcium chloride binds with water and solidifies, it forms very rare minerals: the dihydrate is called sinjarite, while the hexahydrate is known as antarcticite. However, it can also be easily ordered online.
        \item \textbf{Ammonia} (10 ml (milliliter)), \textit{Components}.
        Ammonia (NH\_3) acts as an essential base chemical in the chemical industry.

Its main area of application is in the production of nitrogenous fertilizers (e.g., urea or ammonium nitrate), which account for the majority of global consumption. In addition, ammonia serves as a precursor for organic syntheses, particularly in the production of polymers and synthetic fibers.

The substance is classified as toxic. In high concentrations, ammonia vapors have a strong irritant effect on the mucous membranes of the eyes and the respiratory system, which can lead to severe lung damage and even death. However, such extreme events are rare in everyday use.
        \item \textbf{Opium Poppy (Papaver somniferum)} (5 Bunch), \textit{Components}.
        The opium poppy (Papaver somniferum) belongs to the poppy family and has a long history as one of our oldest medicinal plants.

However, it is not only of interest for medicinal purposes: its seeds are valued in cooking as a foodstuff or pressed to produce high-quality edible oil.

What is special about the opium poppy is that almost all parts of the plant contain alkaloids such as morphine. These active ingredients are concentrated mainly in the white milky sap that runs through the entire plant via a fine channel system – this network is particularly dense in the wall of the capsule fruit. To obtain these substances, the unripe capsules are lightly scored so that the sap can escape. When this milky sap dries in the air, it produces the well-known narcotic opium. Incidentally, the name is derived from Greek and simply means “little sap.”
    \end{denseitemize}


    \subsection*{Venom of a flying snake}

    \begin{tabular}{@{}ll@{}}
        \textbf{\trans{difficulty_label}:} & simple \\
    \end{tabular}

    Preparation:
First, carefully heat the raw venom of the flying snake in a small bowl. But beware of boiling heat! When it begins to steam, the volatile paralyzing substances evaporate. Add the crushed celandine to stimulate blood circulation at the bite site on the victim, allowing the poison to enter the body's fluids more quickly. Stir in the dried sage until the liquid takes on a cloudy, amber color. Finally, press the mixture through a fine linen cloth or a filter made of charcoal dust to remove the coarse residues. What remains is a clear, slightly viscous liquid.

Preparation time: 30 minutes of active work, followed by about 1 hour of resting time to cool and settle.

Quantity: 1 application (approx. 5 Oth of liquid), total weight approx. 5 Oth.

Effect:
The victim feels an immediate coldness at the wound, which spreads like lead into the limbs. It is not fatal, but leads to severe dizziness and slowed reflexes.
The victim must pass a resistance check against MW 8. If they fail, their movement range is reduced by 2 points for 1D6 combat rounds. The victim also suffers 1D3 points of exhaustion.

    \paragraph{\trans{ingredients_label}}
    \begin{denseitemize}
        \item \textbf{Salvia (Salvia officinalis)} (1 Bunch), \textit{Components}.
        The leaves of salvia have an anti-inflammatory, antiperspirant and astringent effect. A tea or rinses are recommended for sore throats or even sweating.
        \item \textbf{Greater celandine (Chelidonium majus)} (1 Bunch), \textit{Components}.
        In the Middle Ages, celandine was used for skin rashes, impaired vision or jaundice. The alkaloids of the plant have an antispasmodic effect. They help with digestive problems and stimulate the flow of bile.
        \item \textbf{Flying Snake Venom Vial} (1 Vial), \textit{Potions and Poisons}.
        A vial filles with the venom of a flying snake.
    \end{denseitemize}

\end{multicols}